\section*{अध्याय \ref{ch:g1:watching-the-weather} --- मौसम को देखना}

\begin{solution}{exo:g1:watching-the-weather:1}
धूप: \omwx{sun}\quad बारिश: \omwx{rain}\quad बर्फ़: \omwx{snow}\quad
बादल: \omwx{cloud}। सूर्य, बूँदों वाले बादल, हिमकणों वाले बादल और अकेले
बादल का कोई भी साफ़ चित्र सही है।
\end{solution}

\begin{solution}{exo:g1:watching-the-weather:2}
पार्क: धूप और गरमी। सड़क: बर्फ़ और ठंड। पहाड़ी: तेज़ हवा (और कुछ बादल)।
स्कूल का फाटक: कोहरा।
\end{solution}

\begin{solution}{exo:g1:watching-the-weather:3}
मंगलवार की सुबह: बादल। शुक्रवार की दोपहर: बौछारें। रविवार की दोपहर:
बारिश।
\end{solution}

\begin{solution}{exo:g1:watching-the-weather:4}
बुधवार: सुबह में भी सूर्य बना है और दोपहर में भी। (गुरुवार की दोपहर में
धूप और बादल हैं।)
\end{solution}

\begin{solution}{exo:g1:watching-the-weather:5}
धूप वाली टोपी गरमी और तेज़ धूप वाले \omterm{def:g1:watching-the-weather:weather}{मौसम} के लिए है; लंबे जूते और बरसाती
बारिश के लिए; गुलूबंद और दस्ताने ठंड और बर्फ़ वाले \omterm{def:g1:watching-the-weather:weather}{मौसम} के लिए।
\end{solution}

\begin{solution}{exo:g1:watching-the-weather:6}
तेज़ हवा वाला \omterm{def:g1:watching-the-weather:weather}{मौसम}: हवा झंडे को धकेल रही है।
\end{solution}

\begin{solution}{exo:g1:watching-the-weather:7}
\begin{enumerate}
  \item ग़लत: दिन भर में \omterm{def:g1:watching-the-weather:weather}{मौसम} बदल सकता है, जैसे धूसर सुबह और धूप वाली
        दोपहर।
  \item सही।
  \item सही: सूर्य इतना चमकीला है कि वह एक पल में भी आँख को हमेशा के लिए
        चोट पहुँचा सकता है।
\end{enumerate}
\end{solution}

\begin{solution}{exo:g1:watching-the-weather:8}
बारिश या बौछारों वाली दोपहरें: सोमवार (बारिश), शुक्रवार (बौछारें) और
रविवार (बारिश), यानी 3। ठंडे दिन: सोमवार, मंगलवार और रविवार, यानी 3।
बाक़ी 4 दिन गरम थे, इसलिए हफ़्ता ज़्यादातर गरम रहा।
\end{solution}

\begin{solution}{exo:g1:watching-the-weather:9}
सुबह: \omwx{cloud} (धूसर आसमान); दोपहर का खाना: \omwx{rain} (बारिश);
शाम: \omwx{suncloud} (कुछ बादलों वाला नीला आसमान)।
\end{solution}

\begin{solution}{exo:g1:watching-the-weather:10}
झंडा उस ओर इशारा करता है जिधर हवा जा रही है। अगर वह बाईं ओर इशारा करता
है, तो हवा बाईं ओर जा रही है, इसलिए वह दाईं ओर से आती है।
\end{solution}

\begin{solution}{exo:g1:watching-the-weather:11}
पत्तियाँ नहीं: सर्दी में। फूल: वसंत में। सर्दी के लिए बर्फ़ का चिह्न
है: सर्दी में ठंड होती है, और कई जगहों पर बर्फ़ सिर्फ़ उसी ऋतु में गिरती है।
\end{solution}
